\chapter{Fortran 77 Basics}

A Fortran program is just a sequence of lines of text. The text has to follow a certain \emph{structure} to be a valid Fortran program. Let's start with the simplest possible program:

\begin{lstlisting}[style=fortran77]
      program hello
c This is a simple "Hello, World!" program in Fortran 77
      
      write (*,*) 'Hello, World!'
      
      stop
      end
\end{lstlisting}

This program contains the minimum elements required in every Fortran program: a program declaration, executable statements (in this case, just one output statement), and an end marker. It simply displays the text "Hello, World!" to the screen.

Now let's look at a slightly more complex example:

\begin{lstlisting}[style=fortran77]
      program circle
      real r, area
 
c This program reads a real number r and prints
c the area of a circle with radius r.
 
      write (*,*) 'Give radius r:'
      read  (*,*) r
      area = 3.14159*r*r
      write (*,*) 'Area = ', area
 
      stop
      end
\end{lstlisting}

The lines that begin with with a "c" are \emph{comments} and have no purpose other than to make the program more readable for humans. Originally, all Fortran programs had to be written in all upper-case letters. Most people now write lower-case since this is more legible, and so will we. You may wish to mix case, but Fortran is not case-sensitive, so "X" and "x" are the same variable.

Here's another example showing basic arithmetic operations:

\begin{lstlisting}[style=fortran77]
      program arithmetic
      integer a, b, sum, diff, prod
      real quotient
      
c This program demonstrates basic arithmetic operations
      
      a = 15
      b = 4
      
      sum = a + b
      diff = a - b
      prod = a * b
      quotient = real(a) / real(b)
      
      write (*,*) 'a = ', a, ' b = ', b
      write (*,*) 'Sum: ', sum
      write (*,*) 'Difference: ', diff
      write (*,*) 'Product: ', prod
      write (*,*) 'Quotient: ', quotient
      
      stop
      end
\end{lstlisting}

Notice in this example how we convert integers to real values before division to ensure we get a decimal result rather than truncated integer division.

\section{Program organization}

A Fortran program generally consists of a main program (or driver) and possibly several subprograms (procedures or subroutines). For now we will place all the statements in the main program; subprograms will be treated later. The structure of a main program is:

\begin{lstlisting}[style=fortran77]
      program name

      declarations

      statements

      stop
      end
\end{lstlisting}

In this tutorial, words that are in \emph{italics} should not be taken as literal text, but rather as a description of what belongs in their place.

The \texttt{stop} statement is optional and may seem superfluous since the program will stop when it reaches the end anyway, but it is recommended to always terminate a program with the \texttt{stop} statement to emphasize that the execution flow stops there.

You should note that you cannot have a variable with the same name as the program.

Here's a practical example showing a temperature conversion program:

\begin{lstlisting}[style=fortran77]
      program temperature
      real fahr, celsius
      
c This program converts Fahrenheit to Celsius
      
      write (*,*) 'Enter temperature in Fahrenheit:'
      read (*,*) fahr
      
      celsius = (fahr - 32.0) * 5.0 / 9.0
      
      write (*,*) fahr, ' degrees Fahrenheit equals'
      write (*,*) celsius, ' degrees Celsius'
      
      stop
      end
\end{lstlisting}

\section{Column position rules}

Fortran 77 is \emph{not} a free-format language, but has a very strict set of rules for how the source code should be formatted. The most important rules are the column position rules:

\begin{lstlisting}[style=fortran77]
Col. 1    : Blank, or a "c" or "*" for comments
Col. 1-5  : Statement label (optional)
Col. 6    : Continuation of previous line (optional)
Col. 7-72 : Statements
Col. 73-80: Sequence number (optional, rarely used today)
\end{lstlisting}

Most lines in a Fortran 77 program starts with 6 blanks and ends before column 72, i.e. only the statement field is used.

\section{Comments}

A line that begins with the letter "c" or an asterisk in the first column is a comment. Comments may appear anywhere in the program. Well-written comments are crucial to program readability. Commercial Fortran codes often contain about 50\% comments. You may also encounter Fortran programs that use the exclamation mark (!) for comments. This is not a standard part of Fortran 77, but is supported by several Fortran 77 compilers and is explicitly allowed in Fortran 90. When understood, the exclamation mark may appear anywhere on a line (except in positions 2-6).

Here's an example demonstrating different comment styles:

\begin{lstlisting}[style=fortran77]
      program comments
c Traditional Fortran 77 comment with 'c' in column 1
* Alternate Fortran 77 comment with '*' in column 1
      
      write (*,*) 'This program demonstrates comment styles'
c     write (*,*) 'This line is commented out'
* The next line shows inline comments (non-standard but commonly supported)
      write (*,*) 'Hello' ! This is an inline comment
      
      stop
      end
\end{lstlisting}

\section{Continuation}

Sometimes, a statement does not fit into the 66 available columns of a single line. One can then break the statement into two or more lines, and use the continuation mark in position 6. Example:

\begin{lstlisting}[style=fortran77]
c23456789 (This demonstrates column position!)

c The next statement goes over two physical lines
      area = 3.14159265358979
     +       * r * r
\end{lstlisting}

Any character can be used instead of the plus sign as a continuation character. It is considered good programming style to use either the plus sign, an ampersand, or digits (using 2 for the second line, 3 for the third, and so on).

Here's a more practical example of using continuation lines for a complex calculation:

\begin{lstlisting}[style=fortran77]
      program long_equation
      real x, y, result
      
c This program calculates a complex formula using continuation lines
      
      x = 2.5
      y = 3.7
      
      result = 3.14159265358979 * x**2 + 2.71828182846 * y**2 
     +         + 1.4142135623 * x * y + (x + y) / 2.0
     +         - (x - y)**2
      
      write (*,*) 'x = ', x, ' y = ', y
      write (*,*) 'Result = ', result
      
      stop
      end
\end{lstlisting}

\section{Blank spaces}

Blank spaces are ignored in Fortran 77. So if you remove all blanks in a Fortran 77 program, the program is still acceptable to a complier but almost unreadable to humans.

\section{Putting it all together}

Let's conclude with a final example that combines several of the concepts we've covered. This program calculates both the area and perimeter of a rectangle:

\begin{lstlisting}[style=fortran77]
      program rectangle
      real length, width, area, perimeter
      
c This program computes the area and perimeter of a rectangle
      
      write (*,*) 'Enter rectangle length:'
      read (*,*) length
      write (*,*) 'Enter rectangle width:'
      read (*,*) width
      
      area = length * width
      perimeter = 2.0 * (length + width)
      
      write (*,*) 'Rectangle dimensions: ', length, ' x ', width
      write (*,*) 'Area: ', area
      write (*,*) 'Perimeter: ', perimeter
      
      stop
      end
\end{lstlisting}

This example demonstrates input/output operations, variable declarations, comments, and arithmetic calculations---all fundamental aspects of Fortran 77 programming.

\section{Exercises}

Solutions are in the directory \texttt{code/exercises/ch03} that accompanies this tutorial.

\begin{enumerate}
    \item Write a program that reads a temperature in degrees Celsius and prints it in degrees Fahrenheit, using $F = \frac{9}{5}C + 32$. Test it with 100 degrees (the answer is 212). Be careful to write \texttt{9.0/5.0} and not \texttt{9/5}; Chapter 6 explains why.
    \item Write a program that reads three numbers and prints their sum and their average.
    \item The following program does not compile, because some of its lines break the column rules. Find the errors and correct them.
\begin{lstlisting}[style=fortran77]
     program errors
      real base, height
      base = 2.0
   height = 3.0
* Area of a triangle
      write (*,*) 'Area of the triangle:',
      0.5*base*height
      end
\end{lstlisting}
    \item Write a program that reads the base and the height of a triangle and prints its area. Start the program with a comment that explains what it does, and split the statement that computes the area over two lines with a continuation line.
    \item Type in the following program, taking care that the number \texttt{25} starts exactly in column 72 (the comment line is a ruler that shows the last digit of each column number). Compile it, run it, and explain the result. Then compile it again with the option \texttt{-Wall}.
\begin{lstlisting}[style=fortran77,basicstyle=\ttfamily\scriptsize]
      program trunc
      real a, b, total
      a = 1.0
      b = 2.0
c23456789012345678901234567890123456789012345678901234567890123456789012
      total = a + b +                                                  25
      write (*,*) 'total =', total
      end
\end{lstlisting}
\end{enumerate}
